\section{Optimiser Comparison on CEC'2005 Functions}
\label{sec:t2-cec}

\subsection{Functions and reproducible instances}
Two functions from the CEC'2005 real-parameter suite~\cite{suganthan2005} provide contrasting
search landscapes. F6 is shifted Rosenbrock, a non-separable problem with a curved valley;
F9 is shifted Rastrigin, a separable problem with many local minima. Their definitions are
\begin{align}
F_6(\mathbf{x}) &= \sum_{i=1}^{D-1}[100(z_i^2-z_{i+1})^2+(z_i-1)^2]+390,\\
F_9(\mathbf{x}) &= \sum_{i=1}^{D}[z_i^2-10\cos(2\pi z_i)+10]-330.
\end{align}
For F6, $\mathbf{z}=\mathbf{x}-\mathbf{o}+1$, with domain $[-100,100]^D$ and optimum 390.
For F9, $\mathbf{z}=\mathbf{x}-\mathbf{o}$, with domain $[-5,5]^D$ and optimum $-330$.
Both use the first $D$ entries of the \emph{official} CEC2005 shift files, matching the reference
MATLAB implementation. The vendored data include source attribution and checksums; earlier
experiments using locally generated shifts are archived and are not the results reported here.
Figure~\ref{fig:t2-landscapes} visualises the two official two-dimensional landscapes.

There are 15 independent repetitions for each algorithm, function and dimension, with
$D\in\{2,10\}$: 180 runs in total. Every run receives $10,000D$ objective evaluations.
The random stream is determined by the repetition, algorithm, dimension and function; streams
are reproducible but do not represent identical initial populations across different algorithms.
An objective wrapper counts evaluations and rejects calls after the budget is exhausted.
The two functions complement each other, but success on them does not establish general
optimiser superiority across the CEC suite.

\subsection{Algorithms and evaluation}
The \textbf{GA} is real-coded, with population 50, tournament size three, blend crossover
($\alpha=0.5$) at probability 0.9, Gaussian mutation with scale 0.1 times the search range and
probability 0.1, and two-member elitism. Real coding avoids discretising these continuous
benchmarks; it differs from the deliberately binary chromosome in the controller experiment.

\textbf{Particle swarm optimisation (PSO)}~\cite{kennedy1995} uses a swarm of 40 particles,
global-best topology, inertia decreasing linearly from 0.9 to 0.4, and cognitive/social
coefficients 1.49445. Velocities are capped at 0.2 times the domain range and positions are
bounded. Its observed poor runs should not be read as a limit on other PSO variants or
attributed to absent velocity control: this implementation already has that control.

\textbf{Simulated annealing (SA)}~\cite{kirkpatrick1983} uses a single candidate with stochastic
proposals and temperature-dependent acceptance of worse moves. Temperature decreases by a
factor of 0.995 after each chain of 20 proposals; the Gaussian proposal scale decreases from
0.20 to 0.002 times the range. Initial temperature is calibrated to accept a typical uphill
move with probability 0.5, and there are no restarts. These settings are stored with the results.
Matching function evaluations
equalises objective-call budgets, not wall-clock time or the suitability of a parameter setting
to each function.

Table~\ref{tab:t2-cec} reports means, sample standard deviations, best and worst final objective
values. Convergence plots show median best-so-far error with an interquartile band, and
Fig.~\ref{fig:t2-cecdist} displays all final-run errors. We also apply two-sided Mann--Whitney
tests with Holm correction across the 12 pairwise comparisons. These test distributional
differences; they are not tests of mean performance, and median differences alone do not
fully describe skewed distributions.

\begin{table*}[t]
\caption{Official CEC2005 instances: final objective values over 15 runs per cell. Lower is better. Standard deviation uses $n-1$. Optima are 390 (F6) and $-330$ (F9).}
\label{tab:t2-cec}\centering\small
\begin{tabular}{@{}llrrrrr@{}}\toprule
Function & $D$ & Algorithm & Mean & Std. dev. & Best & Worst\\\midrule
F6 & 2 & GA & 396.728071 & 20.776039 & 390.000003 & 469.881943\\
F6 & 2 & PSO & 390.000000 & 0.000000 & 390.000000 & 390.000000\\
F6 & 2 & SA & 396.144878 & 6.539262 & 390.305091 & 410.602213\\
\midrule
F6 & 10 & GA & 465.166890 & 86.859403 & 395.757002 & 715.564526\\
F6 & 10 & PSO & $4.031\times10^{5}$ & $8.337\times10^{5}$ & 390.237147 & $2.014\times10^{6}$\\
F6 & 10 & SA & 1506.071423 & 2048.856085 & 428.638385 & 7670.107816\\
\midrule
F9 & 2 & GA & -330.000000 & 0.000000 & -330.000000 & -330.000000\\
F9 & 2 & PSO & -330.000000 & 0.000000 & -330.000000 & -330.000000\\
F9 & 2 & SA & -329.976837 & 0.047067 & -329.999955 & -329.844730\\
\midrule
F9 & 10 & GA & -329.992401 & 0.005209 & -329.997721 & -329.981570\\
F9 & 10 & PSO & -325.031999 & 5.041317 & -329.005041 & -311.751448\\
F9 & 10 & SA & -310.443145 & 10.497144 & -322.738317 & -288.878826\\
\bottomrule\end{tabular}\end{table*}


\subsection{Results and interpretation}
On \textbf{F6 at $D=2$}, PSO reaches the optimum to numerical precision and has the smallest
mean error. GA's mean value is 396.728, slightly worse than SA's 396.145, because one poor
GA run raises its mean. Yet GA's median is 390.025 versus SA's 394.243, and the distributions
differ after correction ($p_{\mathrm{Holm}}=0.0045$). Mean and typical-run comparisons therefore
answer different questions even in two dimensions. PSO differs from both alternatives after
correction ($p_{\mathrm{Holm}}<0.001$).

On \textbf{F9 at $D=2$}, both GA and PSO reach values extremely close to $-330$. Their
corrected comparison does not meet the 0.05 threshold ($p_{\mathrm{Holm}}=0.0559$), and their
practical difference is negligible at the displayed precision. Both improve substantially on
SA's typical final error, although SA also gets close to the optimum in some runs.

On \textbf{F6 at $D=10$}, the ranking is particularly sensitive to summary choice. GA has the
lowest mean objective, 465.167, while PSO's mean is approximately $4.03\times10^5$. However,
PSO's median is 391.661, lower than GA's 468.858. Poor PSO runs dominate its mean and spread.
The GA--PSO distributional comparison is not significant after correcting all 12 tests
($p_{\mathrm{Holm}}=0.0559$). This does not establish equivalence: it shows why a claim that
GA is unconditionally superior would exceed the evidence. SA's median is 750.341, and its
distribution differs from both GA and PSO after correction.

On \textbf{F9 at $D=10$}, GA is best on both mean and median error: mean error above the
optimum is 0.00760, compared with 4.968 for PSO and 19.557 for SA. All three pairwise
comparisons remain significant after correction. The advantage is clearer here than on
Rosenbrock because it is not explained by a few extreme failures in one competitor.

\begin{figure*}[t]\centering
\includegraphics[width=.85\textwidth]{fig14_convergence_cec2005.png}
\caption{Official CEC2005 instances: median best-so-far error and interquartile band over
15 runs. Each cell uses the same evaluation budget across algorithms. Errors below $10^{-12}$
are displayed at that floor for the logarithmic axis.}
\label{fig:t2-cecconv}\end{figure*}

\begin{figure*}[t]\centering
\includegraphics[width=.98\textwidth]{fig18_benchmark_distributions.png}
\caption{All 180 final-run errors with per-cell medians. The large PSO tail on F6 at $D=10$
explains the disagreement between mean and median rankings. A $10^{-12}$ plotting floor is
used for errors at numerical precision.}
\label{fig:t2-cecdist}\end{figure*}

These observations concern the implemented algorithms and fixed parameter settings. Tuning
the annealing schedule, changing PSO's velocity cap or adding restarts, or changing GA variation could alter the
outcomes and would constitute a new experiment. The benchmark does not reproduce the
membership-function constraints or inference cost of the controller; its rankings cannot simply
be transferred to the assistive-care tuning problem. Runtime also differs across algorithms,
so equal evaluation budgets should be interpreted as search-effort control, not equal elapsed time.

\section{Conclusions}
\label{sec:t2-conclusion}
The Mamdani controller provides a demonstrable mapping from room conditions and user
preferences to heating/cooling and lighting actions. Membership functions, a linguistic rule
base, rule-activation examples and response surfaces make the design inspectable. The
reference controller and MATLAB implementation support reproducible examination of the
chosen scenario; simulation alone does not establish effectiveness for real residents.

Genetic tuning reduces normalised test error by 32.6\% relative to the hand-designed baseline
on the synthetic reference-policy dataset. It also exposes a weakness of error-only optimisation:
three of five tuned controllers reorder linguistic sets, four leave some inputs without a defined
output, and none preserves both audited properties. The lower-error controller should therefore
be treated as an experimental result rather than a deployment-ready improvement. A stronger
next design would enforce coverage and ordering as feasibility constraints, then compare error
among valid controllers using the same training/test split.

The official CEC benchmark comparison reinforces the need to inspect distributions. PSO is
strong in two dimensions; GA gives the lowest mean objective in both ten-dimensional cells,
but the F6 median and tail tell a different story. Corrected rank tests qualify claims that
would appear decisive from means alone. No universal winner follows from two functions,
two dimensions and 15 repetitions.

The central limitations are the synthetic target policy, absence of real-resident evaluation,
small number of controller-tuning runs and fixed optimiser parameters. The work demonstrates
the requested design, optimisation and comparison processes while making the consequences
of their assumptions visible.

\section*{Individual Contributions}
Bibek Paudyal (250288) coordinated the fuzzy-controller design and Python implementation.
Sajan Mahat (250289) coordinated genetic tuning, the matched control and the coverage and
linguistic-ordering audits. Siddhartha Bhatta (250620) coordinated the MATLAB implementation
and benchmark comparison. All three members contributed equally to report preparation and
review of the implementation, experiments and conclusions, with shared responsibility for
the final submission.
